\section{Background}
\label{sec:background}

This section establishes the mathematical foundations for our work, 
reviewing classical results on magic squares and introducing the geometric and statistical concepts 
that underpin the Magic Gem framework.

\subsection{Magic Squares: Classical Results}

The enumeration of magic squares reveals a rapid growth in complexity with the order $n$.
For $n = 3$, there exist exactly eight magic squares, all related by rotations and reflections,
yielding a single essentially different magic square---the Lo Shu.
The situation changes dramatically for larger orders: 
$n = 4$ admits 7,040 magic squares forming 880 equivalence classes under the dihedral group,
while $n = 5$ yields approximately 275 million magic squares comprising roughly 34 million equivalence classes.
This explosive growth poses significant challenges for enumeration and classification,
motivating the search for structural insights that transcend explicit construction.

The magic constant $M(n) = n(n^2 + 1)/2$ admits a simple derivation from first principles.
The sum of all entries equals $\sum_{k=1}^{n^2} k = n^2(n^2 + 1)/2$,
which must be distributed equally among the $n$ rows, yielding $M(n)$ per row.
The same argument applies to columns, and the diagonal constraints provide additional structure
that severely restricts the space of valid configurations.

Classical construction methods provide algorithmic approaches to generating magic squares of various orders.
The Siamese method, attributed to de la Loubère, constructs magic squares for odd $n$
by starting in the middle of the top row and moving diagonally up-right,
wrapping around edges and moving down when blocked.
For doubly-even orders ($n$ divisible by 4), one fills the grid sequentially 
and then swaps diagonal elements with their complements.
Singly-even orders ($n = 4k + 2$) require more sophisticated techniques
combining quadrant-based construction with targeted transpositions.
These methods demonstrate that magic squares exist for all $n \geq 3$,
though they produce only a tiny fraction of all possible magic squares for larger orders.

\subsection{Symmetry and the Dihedral Group}

The dihedral group $D_4$ acts naturally on $n \times n$ grids through rotations and reflections.
This group comprises eight elements: four rotations (by $0^\circ$, $90^\circ$, $180^\circ$, and $270^\circ$) 
and four reflections (horizontal, vertical, and along both diagonals).
Since the magic square conditions treat rows, columns, and diagonals symmetrically 
under these operations, $D_4$ permutes magic squares among themselves.

\begin{proposition}
If $S$ is a magic square, then $g \cdot S$ is also a magic square for any $g \in D_4$.
\end{proposition}

\begin{proof}
The $D_4$ action preserves the set of row sums, column sums, and diagonal sums,
merely permuting which lines play which roles.
Since the action also preserves the set of entries $\{1, \ldots, n^2\}$,
it maps magic squares to magic squares.
\end{proof}

This symmetry motivates the notion of \emph{essentially different} magic squares:
two squares are essentially the same if one can be obtained from the other by a $D_4$ transformation.
For $n = 3$, all eight magic squares form a single equivalence class,
while larger orders admit multiple classes with distinct structural properties.

\subsection{Convex Hulls and Polyhedra}

The convex hull of a finite point set $P \subset \R^d$ is the smallest convex set containing $P$,
equivalently characterized as the set of all convex combinations of points in $P$:
\[
\text{conv}(P) = \left\{ \sum_{i=1}^{k} \lambda_i p_i : p_i \in P, \lambda_i \geq 0, \sum_i \lambda_i = 1 \right\}.
\]
For finite point sets in $\R^3$, the convex hull is a convex polyhedron 
whose vertices form a subset of the original points.
The remaining points lie in the interior or on faces of the polyhedron.

Key geometric quantities associated with a convex polyhedron include its volume, 
surface area, and the number of vertices, edges, and faces.
These invariants provide a coarse characterization of the polyhedron's shape
and will serve as descriptors for comparing Magic Gems across different arrangements.

\subsection{Covariance and the Moment of Inertia}

For a discrete collection of points $\{(\mathbf{x}_i)\}_{i=1}^{N}$ with positions $\mathbf{x}_i \in \R^d$,
the covariance matrix captures the second-order statistical structure:
\begin{equation}
\Sigma_{jk} = \frac{1}{N} \sum_{i=1}^{N} (x_{i,j} - \bar{x}_j)(x_{i,k} - \bar{x}_k),
\end{equation}
where $\bar{x}_j = \frac{1}{N} \sum_i x_{i,j}$ denotes the mean of the $j$-th coordinate.
When the point cloud is centered so that $\bar{\mathbf{x}} = \mathbf{0}$, 
this simplifies to $\Sigma_{jk} = \frac{1}{N} \sum_i x_{i,j} x_{i,k}$.

The diagonal entries $\Sigma_{jj}$ measure the variance along each coordinate axis,
while the off-diagonal entries $\Sigma_{jk}$ ($j \neq k$) measure the covariance between coordinates.
A zero off-diagonal entry indicates that the corresponding coordinates are uncorrelated:
knowing one provides no linear information about the other.

Closely related is the moment of inertia tensor from classical mechanics.
For point masses $\{(m_i, \mathbf{x}_i)\}$ with positions $\mathbf{x}_i \in \R^3$, 
the inertia tensor is given by
\begin{equation}
I_{jk} = \sum_{i=1}^{N} m_i \left( \|\mathbf{x}_i\|^2 \delta_{jk} - x_{i,j} x_{i,k} \right).
\end{equation}
The eigenvalues of $I$, called the principal moments of inertia, 
characterize the rotational properties of the point configuration
and are invariant under rigid rotations.
For equal unit masses, the inertia tensor and covariance matrix are closely related,
differing by a trace term and sign conventions.
